\section{Сложность и точность}
\label{sec:complexity}

Обозначения: $n$ --- число вершин, $m$ --- число рёбер, $P$ --- число потоков, $T$ --- число итераций
до сходимости, $B$ --- размер попоточного буфера.

\subsection{Время}
\paragraph{Предобработка (один раз).} Проход по входу с подсчётом степеней и переименованием ---
$\bigO(m)$. Префиксная сумма для \texttt{sourcesPtr} --- $\bigO(n)$. Раздача рёбер по корзинам ---
$\bigO(m)$. Сортировка корзин --- около $\bigO(m\log(m/K))$ при $K$ корзинах, с малым множителем
(сортируются короткие корзины). Итого предобработка --- $\bigO(m\log m)$ по времени и $\Theta(m)$ по
вводу-выводу (несколько последовательных проходов).

\paragraph{Одна итерация.} Вычисление $\mathrm{contrib}$, инициализация $s_{\mathrm{new}}$, суммы ---
$\bigO(n)$. Ядро~\eqref{eq:pull} --- ровно один проход по файлу \texttt{sources}, $\bigO(m)$ сложений
и $\Theta(m)$ последовательного чтения. Распараллеливание делит работу на $P$ потоков. При
сбалансированных по числу рёбер кусках счётное время итерации --- $\bigO((n+m)/P)$, а время чтения
--- $\Theta(m/\beta)$, где $\beta$ --- эффективная скорость диска. Полное время:
\[
  T_{\text{итер}} = \bigO\!\Big(T\cdot\tfrac{n+m}{P}\Big)\ \text{(счёт)}
    \;+\; \Theta\!\Big(T\cdot\tfrac{m}{\beta}\Big)\ \text{(ввод-вывод)}.
\]
Реальное время --- максимум из двух слагаемых. На больших графах доминирует ввод-вывод ($T$
перечитываний \texttt{sources}), поэтому число читаемых байт сведено к минимуму, а чтение идёт
последовательно.

\subsection{Память}
ОЗУ: $\Theta(n)$ на общие векторы $+\ \bigO(P\cdot B)$ на буферы $+\ \bigO(1)$ накладных JVM. Главное
--- не зависит от $m$ и от наличия гипер-узлов. Диск: $\Theta(m)$ на \texttt{sources} $+\ \Theta(n)$
на остальное.

\subsection{Балансировка: средний и худший случай}
\begin{itemize}[itemsep=2pt]
  \item \textbf{Средний случай.} Куски выровнены по числу рёбер, нагрузка распределена ровно,
        ускорение близко к линейному до потолка ввода-вывода.
  \item \textbf{Худший случай (гипер-узел).} Одна вершина-приёмник степени $D$ неделима, поэтому
        отставание ограничено величиной $\bigO(D)$. При разбиении на $K\gg P$ кусков и $D\ll m/K$ это
        отставание незаметно. Патология наступила бы, только если $D$ сравнима с $m/P$, то есть одна
        вершина собирает заметную долю всех рёбер.
\end{itemize}

\subsection{Сходимость}
Степенной метод на расширенном графе сходится к единственному пределу. Двусторонний фоновый узел
делает цепь Маркова сильно связной и апериодичной, отсюда (теорема Перрона--Фробениуса) существование
и единственность стационарного распределения. Заодно это решает проблему висячих вершин: их масса
уходит фоновому узлу и потом перераспределяется. Ошибка убывает геометрически:
$\|s^{(t)} - s^\ast\|_1 \le C\,\rho^{\,t}$, где $\rho<1$ определяется вторым по модулю собственным
значением матрицы перехода. Число итераций $T = \bigO\!\big(\log(1/\varepsilon)/\log(1/\rho)\big)$.

\subsection{Точность и воспроизводимость}
Все суммы (и в ячейках~\eqref{eq:pull}, и в редукциях) считаются в фиксированном порядке, заданном
раскладкой файла \texttt{sources} и порядком потоков. Сложение \texttt{double} не ассоциативно,
поэтому именно фиксированный порядок даёт побитовую воспроизводимость и независимость от числа
потоков. Накопление ошибки округления контролируется редкой перенормировкой к сохраняемой сумме масс
$n$. Про конкретный критерий остановки и наблюдаемое число итераций на реальных графах --- в
разделе~\ref{sec:testing}.
